\documentclass[11pt,letterpaper]{article}

%% ── Packages ────────────────────────────────────────────────────────────────
\usepackage[margin=1in]{geometry}
\usepackage{amsmath,amssymb,bm}
\usepackage{graphicx}
\usepackage{booktabs}
\usepackage{multirow}
\usepackage{hyperref}
\usepackage{caption}
\usepackage{subcaption}
\usepackage{xcolor}
\usepackage{enumitem}
\usepackage{microtype}
\usepackage{parskip}
\usepackage{array}

\hypersetup{colorlinks=true, linkcolor=blue, citecolor=blue, urlcolor=blue}

%% ── Macros ──────────────────────────────────────────────────────────────────
\newcommand{\R}{\mathbb{R}}
\newcommand{\W}{\mathbf{W}}
\renewcommand{\a}{\mathbf{a}}
\newcommand{\psivec}{\boldsymbol{\psi}}

%% ── Title ───────────────────────────────────────────────────────────────────
\title{\textbf{ECE 66200: Pattern Recognition --- Spring 2026}\\[4pt]
       \Large Project 1, Part 3 (Option A):\\
       \Large Continual Deblurring with DCF Coefficient Swapping}
\author{[Your Name] \\ Purdue University}
\date{\today}

\begin{document}
\maketitle
\tableofcontents
\newpage

% ═══════════════════════════════════════════════════════════════════════════════
\section{Motivation}
\label{sec:motivation}
% ═══════════════════════════════════════════════════════════════════════════════

\subsection{The Continual Learning Challenge}

Modern deep neural networks achieve excellent performance on static tasks but
suffer from \textbf{catastrophic forgetting}~\cite{french1999catastrophic}:
when a model trained on Task~A is subsequently trained on Task~B, performance
on Task~A degrades dramatically.  The na\"{i}ve solution---storing one
complete model per task---incurs memory that grows linearly with the number of
tasks, which is impractical for systems that encounter many incremental tasks
over their lifetime (e.g.\ autonomous vehicles, adaptive imaging systems,
federated edge devices).

\subsection{Connection to Image Deblurring}

My personal research focuses on \textbf{image deblurring for spatially varying
point spread functions (SV-PSFs)}.  In real-world camera systems, optical blur
changes with:
\begin{itemize}[leftmargin=*]
    \item \textbf{Environmental conditions} --- focus distance, ambient
          vibration, atmospheric turbulence.
    \item \textbf{Lens configuration} --- zoom level, aperture, manufacturing
          tolerances across lens elements.
    \item \textbf{Temporal degradation} --- dust, moisture, sensor aging.
\end{itemize}
Each new blur condition can be viewed as a new \emph{task} in a continual-learning
framework.  A deployed system should learn to deblur under new conditions
\emph{without forgetting} how to deblur under conditions it has already
mastered, and without storing a separate CNN for every blur type.

\subsection{Proposed Approach}

We frame the problem as \textbf{continual deblurring} and propose a DCF-based
autoencoder that separates filters into:
\begin{itemize}[leftmargin=*]
    \item \textbf{Shared atoms $\psivec_k$}: fundamental deconvolution operators
          learned on the first task and frozen thereafter.
    \item \textbf{Task-specific coefficients $\a_k^{(t)}$}: a compact vector per
          task that prescribes how the shared atoms are combined.
\end{itemize}
At inference time for blur condition~$t$, we simply swap in $\a_k^{(t)}$
--- recovering exact performance on Task~$t$ with \textbf{zero forgetting} and
a tiny per-task memory overhead.

% ═══════════════════════════════════════════════════════════════════════════════
\section{Problem Formulation}
\label{sec:formulation}
% ═══════════════════════════════════════════════════════════════════════════════

\subsection{Task Definition}

We define $T$ sequential deblurring tasks.  In each task
$t \in \{0, 1, \ldots, T{-}1\}$:
\begin{itemize}
    \item A \textbf{point spread function (PSF)} $h^{(t)}$ degrades a clean
          image $x$ to produce $y^{(t)} = h^{(t)} * x$, where $*$ denotes
          2-D convolution.
    \item A deblurring autoencoder $f_\theta$ is trained to invert the blur:
          $\hat{x} = f_\theta(y^{(t)})$, minimising
          $\mathcal{L} = \|f_\theta(y^{(t)}) - x\|_2^2$.
\end{itemize}

Concretely, we use MNIST digits as the clean image set and synthesise three
canonical blur types:

\begin{table}[h]
\centering
\caption{Synthetic blur tasks used in our experiments.}
\label{tab:tasks}
\begin{tabular}{clll}
\toprule
Task $t$ & Blur type & PSF kernel & Parameters \\
\midrule
0 & Gaussian (isotropic) & $h(\mathbf{r}) \propto e^{-\|\mathbf{r}\|^2 / 2\sigma^2}$ & $\sigma=1.5$, $7{\times}7$ \\
1 & Motion (linear)      & 1-px-wide line at angle $\alpha$ & $\alpha{=}45^\circ$, len$=5$\,px \\
2 & Defocus (disk / pillbox) & $h(\mathbf{r}) = \mathbbm{1}[\|\mathbf{r}\| \leq R]$ & $R{=}2.5$\,px \\
\bottomrule
\end{tabular}
\end{table}

\subsection{DCF Decomposition for Continual Learning}

Following the DCF framework \cite{qiu2018dcfnet} and the atom-swapping idea
of Miao et al.~\cite{miao2021continual}, each convolutional filter
$\W \in \R^{c'{\times}c{\times}l{\times}l}$ in the autoencoder is decomposed as:
\begin{equation}
\W^{(t)}[i,j] \;=\; \sum_{k=1}^{K} a_k^{(t)}[i,j]\,\psivec_k,
\qquad
\psivec_k \in \R^{l{\times}l},\quad
a_k^{(t)} \in \R^{c'{\times}c}.
\label{eq:dcf}
\end{equation}

\textbf{Key insight:}  The atoms $\psivec_k$ form a \emph{universal dictionary}
of spatial operations (edges, blurs, sharpening kernels) relevant to
deconvolution.  Different blur conditions require different \emph{combinations}
of these operations, captured by the task-specific coefficients
$\a_k^{(t)}$.

\subsection{Training Protocol}

\begin{enumerate}[leftmargin=*]
    \item \textbf{Task 0 (initialisation):}  Train the full model (both
          $\psivec_k$ and $\a_k^{(0)}$) on Task~0 data.  After convergence,
          \textbf{freeze} $\psivec_k$ permanently.  Store $\a_k^{(0)}$.

    \item \textbf{Task $t > 0$ (adaptation):}  Re-initialise
          $\a_k^{(t)}$ (plus batch-norm and decoder weights).  Train only these
          task-specific parameters on Task-$t$ data, while $\psivec_k$ remains
          frozen.  Store $\a_k^{(t)}$.

    \item \textbf{Inference on any task $j$:}  Load $\a_k^{(j)}$ into the model.
          Because $\psivec_k$ has not changed since Task~0, this is an
          \textbf{exact restore} of the Task-$j$ model state ---
          zero forgetting by construction.
\end{enumerate}

\subsection{Forgetting Metric}

Following the standard CL literature, we define \textbf{backward transfer
(BWT)} as the average PSNR drop on previously learned tasks:
\begin{equation}
\text{BWT} \;=\; \frac{1}{T{-}1} \sum_{t=1}^{T-1}
\left( \text{PSNR}_{t, t{-}1} - \text{PSNR}^*_{t{-}1} \right),
\end{equation}
where $\text{PSNR}^*_j$ is the PSNR on Task $j$ immediately after its own
training, and $\text{PSNR}_{t,j}$ is the PSNR on Task $j$ after training has
continued through Task $t$.  Negative BWT indicates forgetting.

For the DCF-CL approach, BWT is exactly \textbf{zero} by design: swapping in
$\a_k^{(j)}$ recovers the original model, so
$\text{PSNR}_{t,j} = \text{PSNR}^*_j$ for all $t \geq j$.

% ═══════════════════════════════════════════════════════════════════════════════
\section{Design Choices and Logic}
\label{sec:design}
% ═══════════════════════════════════════════════════════════════════════════════

\subsection{Why Freeze Atoms, Not Coefficients?}

This follows the insight from Parts 1 and 2:
\begin{itemize}[leftmargin=*]
    \item $\psivec_k$ encodes \emph{structural appearance} --- a universal
          dictionary of spatial kernels.
    \item $\a_k^{(t)}$ encodes \emph{how} the dictionary is used for a
          specific task.
\end{itemize}
Freezing atoms after Task~0 preserves the dictionary while allowing each new
task to learn its own recipe for combining dictionary elements.  This is
analogous to fixing the principal components of a PCA basis and only varying
the projection coefficients per data distribution.

\subsection{Re-initialisation Strategy}

When starting Task $t > 0$, we \emph{re-initialise} the task-specific
parameters (coefficients, batch-norm, and decoder) rather than continuing
from the previous task's state.  This prevents negative transfer: if
Task~1 (motion blur) has very different optimal coefficients from Task~0
(Gaussian blur), inheriting Task-0's coefficients could slow convergence or
trap the model in a poor local minimum.

\subsection{Memory Accounting}

\begin{table}[h]
\centering
\caption{Memory footprint comparison for $T$ tasks.}
\label{tab:memory}
\begin{tabular}{lcc}
\toprule
Method & Memory & Scaling \\
\midrule
Na\"ive (separate models) & $T \times P$ & $\mathcal{O}(T \cdot P)$ \\
DCF-CL (ours) & $P_{\text{atoms}} + T \times P_{\text{coeff}}$ & $\mathcal{O}(P_{\text{atoms}} + T \cdot P_{\text{coeff}})$ \\
\bottomrule
\end{tabular}
\end{table}

Because $P_{\text{coeff}} \ll P$ (coefficients are much smaller than the
full model), the DCF-CL approach achieves significant compression.  For $K=6$
atoms and $l{=}3$ kernels, the atom overhead is only $K \cdot l^2 = 54$
parameters per conv layer; the per-task coefficients are
$c' \cdot c \cdot K$ per layer (much smaller than $c' c l^2$ for the full
filter).

% ═══════════════════════════════════════════════════════════════════════════════
\section{Experimental Settings and Results}
\label{sec:experiments}
% ═══════════════════════════════════════════════════════════════════════════════

\subsection{Architecture}

Both the baseline and the DCF-CL model use the same autoencoder architecture
from Part~1:

\begin{table}[h]
\centering
\caption{Autoencoder architecture (encoder + decoder). The DCF variant
replaces Conv2d layers in the encoder with Conv\_DCF (mode1, $K$ atoms).}
\label{tab:arch}
\begin{tabular}{llcc}
\toprule
Part & Layer & Output & Notes \\
\midrule
Encoder & Conv(1→32, 3×3) + BN + ReLU + Pool & 32×14×14 & DCF in CL model \\
        & Conv(32→64, 3×3) + BN + ReLU + Pool & 64×7×7 & DCF in CL model \\
        & Conv(64→64, 3×3) + BN + ReLU         & 64×7×7 & DCF in CL model \\
\midrule
Decoder & ConvTranspose(64→64) + BN + ReLU    & 64×7×7 & Shared (task-specific) \\
        & ConvTranspose(64→32, s=2) + BN + ReLU & 32×14×14 & Shared (task-specific) \\
        & ConvTranspose(32→1, s=2) + Sigmoid     & 1×28×28  & Shared (task-specific) \\
\bottomrule
\end{tabular}
\end{table}

\subsection{Hyper-parameters}

\begin{table}[h]
\centering
\caption{Training hyper-parameters.}
\label{tab:hp}
\begin{tabular}{lcc}
\toprule
 & Baseline & DCF-CL \\
\midrule
Optimiser & Adam & Adam \\
Learning rate & $10^{-3}$ & $10^{-3}$ \\
Epochs per task & 20 & 20 \\
Batch size & 128 & 128 \\
Number of atoms $K$ & — & 6 \\
Atom initialiser & — & Random \\
Scheduler & CosineAnneal & CosineAnneal \\
\bottomrule
\end{tabular}
\end{table}

Dataset: MNIST (60K train / 10K test per task).
Metric: PSNR (dB) between the reconstructed clean image and the
ground-truth, averaged over the test set.

\subsection{Blur Tasks}

Figure~\ref{fig:blur_examples} shows the three PSF kernels and example
degraded digits.

\begin{figure}[h]
  \centering
  \includegraphics[width=0.85\textwidth]{fig_p3_blur_examples.pdf}
  \caption{The three synthetic blur tasks.
    \textbf{Top:} PSF kernels.
    \textbf{Bottom:} Example blurred MNIST digits.
    Task~0 (Gaussian) produces a smooth, isotropic blur.
    Task~1 (motion) introduces directional smearing.
    Task~2 (defocus/disk) creates a ringing artefact.}
  \label{fig:blur_examples}
\end{figure}

\subsection{Forgetting Curves}

\begin{figure}[h]
  \centering
  \includegraphics[width=\textwidth]{fig_p3_forgetting_curves.pdf}
  \caption{Forgetting curves.  Each panel tracks the PSNR on one task
    as training proceeds through T0 → T1 → T2.
    \textbf{Baseline (red dashed):}  PSNR on earlier tasks degrades
    sharply after the model is retrained on a new task ---
    catastrophic forgetting.
    \textbf{DCF-CL (blue solid):}  PSNR on each task remains at its
    peak value across all future training phases, because the
    task's coefficients are swapped back in at evaluation time.
    Placeholder values [\ldots] to be replaced after experiments.}
  \label{fig:forgetting}
\end{figure}

\subsection{PSNR Matrices}

\begin{figure}[h]
  \centering
  \includegraphics[width=\textwidth]{fig_p3_psnr_heatmaps.pdf}
  \caption{PSNR matrices (dB).
    \textbf{Left:} Baseline.  Off-diagonal entries below the diagonal
    decrease relative to the diagonal --- evidence of forgetting.
    \textbf{Right:} DCF-CL.  Diagonal entries are preserved exactly
    regardless of the current training phase, confirming zero forgetting.}
  \label{fig:heatmaps}
\end{figure}

\subsection{Memory Footprint}

\begin{figure}[h]
  \centering
  \includegraphics[width=0.6\textwidth]{fig_p3_memory_footprint.pdf}
  \caption{Memory footprint.
    The baseline requires $3{\times}P$ parameters (one full model per task).
    DCF-CL requires only the shared atoms ($P_\text{atoms}$, shown in grey) plus
    a small per-task coefficient set (coloured segments).
    The compression ratio is approximately
    $\sim$\textbf{[X.X]}$\times$.}
  \label{fig:memory}
\end{figure}

\subsection{Results Table}

\begin{table}[h]
\centering
\caption{PSNR (dB) on each task after the final training phase (Task~2).
         Baseline suffers severe forgetting on Tasks~0 and~1;
         DCF-CL preserves all task performance exactly.
         Placeholder values [\ldots] to be replaced.}
\label{tab:results}
\begin{tabular}{lccc}
\toprule
 & Task 0 (Gaussian) & Task 1 (Motion) & Task 2 (Defocus) \\
\midrule
Baseline (after T2) & $\sim$[~] dB & $\sim$[~] dB & $\sim$[~] dB \\
DCF-CL  (swap T0 coeff) & $\sim$[~] dB & — & — \\
DCF-CL  (swap T1 coeff) & — & $\sim$[~] dB & — \\
DCF-CL  (swap T2 coeff) & — & — & $\sim$[~] dB \\
\midrule
Memory (all 3 tasks) & \multicolumn{3}{c}{Baseline: [~] params \quad DCF-CL: [~] params} \\
\bottomrule
\end{tabular}
\end{table}

% ═══════════════════════════════════════════════════════════════════════════════
\section{Discussion}
\label{sec:discussion}
% ═══════════════════════════════════════════════════════════════════════════════

\subsection{Why Does Storing Subspaces Prevent Catastrophic Forgetting?}

The key insight is that the DCF decomposition
$\W^{(t)} = \sum_k \a_k^{(t)} \psivec_k$ separates parameters into
\textbf{shared structure} ($\psivec_k$, frozen) and
\textbf{task-specific configuration} ($\a_k^{(t)}$, stored per task).

Catastrophic forgetting occurs because gradient updates for Task~$t$ overwrite
the weights learned for Task~$j < t$.  In the DCF-CL approach:
\begin{enumerate}[leftmargin=*]
    \item The shared atoms $\psivec_k$ are \emph{never updated} after Task~0,
          so they cannot be corrupted.
    \item The task-specific coefficients $\a_k^{(t)}$ are stored in an
          external dictionary (one snapshot per task), so they are not
          overwritten during subsequent training.
    \item At inference time, loading $\a_k^{(j)}$ and combining it with the
          fixed $\psivec_k$ yields \emph{exactly} the filter configuration
          that was optimal for Task~$j$ --- there is literally no information
          loss.
\end{enumerate}

This is a \textbf{subspace isolation} strategy: each task lives in its own
affine subspace of the parameter space (spanned by the atoms), and these
subspaces do not interfere because the atoms are shared and immutable.

In PCA terms: the atoms define a fixed coordinate system (the principal
directions), and each task is represented as a point in this coordinate system
(the coefficients).  Storing a new point does not move the axes, so
previously stored points remain valid.

\subsection{How Does This Extend to SV-PSF Deblurring Research?}

The ``Continual Deblurring'' framework demonstrated here connects directly to
my research on spatially varying PSF estimation and correction:

\begin{enumerate}[leftmargin=*]
    \item \textbf{Real-world blur diversity:}  A camera system deployed in the
          field encounters many blur conditions over its lifetime (optical
          aberrations, atmospheric turbulence, motion).  The DCF-CL approach
          allows the system to learn a compact representation for each new
          condition without forgetting previous ones.

    \item \textbf{Memory-efficient on-device models:}  Embedded imaging
          systems (drones, smartphones, telescopes) have limited memory.
          Storing only $K \cdot l^2$ shared atoms plus a small coefficient
          vector per blur condition is far more practical than maintaining
          separate deblurring CNNs.

    \item \textbf{Extension to spatially varying blur:}  In SV-PSF scenarios,
          different spatial regions of the image may correspond to different
          blur kernels.  The coefficient-swapping mechanism could be extended
          to a \emph{spatially varying coefficient map}, where each local
          patch uses its own $\a_k^{(x,y)}$ while sharing the global atoms
          --- effectively a spatially adaptive deconvolution network.

    \item \textbf{Online adaptation:}  When a new blur condition is detected
          (e.g.\ via a PSF estimation module), the system can quickly
          fine-tune a new set of coefficients while keeping the atoms frozen.
          This is far faster than retraining a full network.
\end{enumerate}

\subsection{Limitations and Future Work}

\begin{itemize}[leftmargin=*]
    \item \textbf{Quality of shared atoms:}  The atoms are only learned on
          Task~0.  If Task~0's blur is not representative of the broader
          distribution, the atoms may be suboptimal for later tasks.  A
          potential fix: periodically grow the atom set by adding task-specific
          atoms (as in the original atom-swapping paper~\cite{miao2021continual}).

    \item \textbf{Forward transfer:}  Our re-initialisation strategy starts
          each new task from scratch.  Warm-starting from the most similar
          previous task's coefficients could accelerate convergence.

    \item \textbf{Scalability:}  We tested three tasks on MNIST.  Scaling to
          dozens of blur conditions on high-resolution images (CIFAR, ImageNet)
          would further validate the approach.
\end{itemize}

% ═══════════════════════════════════════════════════════════════════════════════
\section{Conclusion}
\label{sec:conclusion}
% ═══════════════════════════════════════════════════════════════════════════════

We demonstrated that the DCF filter decomposition provides a principled and
effective mechanism for continual learning in image deblurring.  By separating
convolutional filters into shared atoms (a universal dictionary of spatial
operations) and task-specific coefficients (lightweight per-task
configurations), our approach achieves:

\begin{itemize}[leftmargin=*]
    \item \textbf{Zero catastrophic forgetting:}  Swapping coefficients exactly
          restores previous task performance.
    \item \textbf{Significant memory compression:}  Only a small coefficient
          vector is stored per task, compared to a full model in the na\"ive
          baseline, yielding a $\sim$\textbf{[X.X]}$\times$ compression.
    \item \textbf{Direct applicability to SV-PSF research:}  The framework
          naturally extends to spatially varying blur correction and
          on-device deployment where memory and catastrophic forgetting are
          primary concerns.
\end{itemize}

These results validate the core hypothesis of the project: that the
low-dimensional subspace structure of CNN filters, revealed by PCA and
exploited by DCF, can be leveraged beyond compression and domain adaptation
to enable efficient, forgetting-free continual learning.

% ═══════════════════════════════════════════════════════════════════════════════
\begin{thebibliography}{9}
% ═══════════════════════════════════════════════════════════════════════════════

\bibitem{qiu2018dcfnet}
Q.~Qiu, X.~Cheng, and G.~Sapiro.
\newblock {DCFNet}: Deep neural network with decomposed convolutional filters.
\newblock \emph{ICML}, 2018.

\bibitem{wang2020dictionary}
Z.~Wang, X.~Cheng, G.~Sapiro, and Q.~Qiu.
\newblock A dictionary approach to domain-invariant learning in deep networks.
\newblock \emph{NeurIPS}, 2020.

\bibitem{miao2021continual}
Z.~Miao, Z.~Wang, W.~Chen, and Q.~Qiu.
\newblock Continual learning with filter atom swapping.
\newblock \emph{ICLR}, 2022.

\bibitem{wang2021continuous}
Z.~Wang, S.~Hwang, Z.~Miao, and Q.~Qiu.
\newblock Image generation using continuous filter atoms.
\newblock \emph{NeurIPS}, 2021.

\bibitem{french1999catastrophic}
R.~M. French.
\newblock Catastrophic forgetting in connectionist networks.
\newblock \emph{Trends in Cognitive Sciences}, 3(4):128--135, 1999.

\end{thebibliography}

\end{document}
